\answer{ex:14:1}{(a)~$0.17$~s y $1.5$~s. (b)~A distancias menores de $11$~cm.}
\answer{ex:14:2}{El tacto sin dolor no pasa con ningún $\mu$ mientras $g\le\beta w_{q\mu}/w_{z\mu}=5$; el tacto $\mu=1$ pasa con $g>6$: una sensibilización fuerte convierte un roce corriente en dolor.}
\answer{ex:14:3}{(a)~$g^*=(1-k_sC)/(1-k_sA)$, estable con $k_sA<1$. (b)~Sin tacto, con $\nu\ge2$; con tacto, ya con $\nu\ge1.7$.}
\answer{ex:14:4}{$V(t)=-Pe^{-(T-t)/\tau_\gamma}$. (a)~$-Pe^{-T/\tau_\gamma}\approx-0.37P$. (b)~$+Pe^{-(T-t_e)/\tau_\gamma}\approx+0.61P$; crece con $t_e$ hasta $+P$ y enseña a evitar el dolor.}
\answer{ex:14:5}{Con $k_s=4$ no hay traba (el segundo estado estable existe con $k_s\ge4.58$); $T_{\min}\approx4.35$, $1.40$, $0.65$, $0.32\,\tau_s$ con $k_s=4.6$, $5$, $6$, $8$.}
\answer{ex:14:6}{$w_{q\nu}=4/9\approx0.444$, $q_0=2/9\approx0.222$, $w_{q\mu}=49/45\approx1.089$; el tacto sin dolor es seguro hasta $g=49/9\approx5.4$.}
\answer{ex:14:7}{(a)~Umbral $\nu_c(g)=\theta/g$ si $\nu_c\ge q_0/w_{q\nu}$; si no, $(\theta+\beta q_0)/(g+\beta w_{q\nu})$; $g^*\approx2.45$, $1.0$, $0.25$. (b)~Pregunta abierta.}
